\documentclass{acmart}

\usepackage{natbib}

\title{Using Alloy to Model Agentic Programming}

\author{Isaac H. Lopez Diaz}

\begin{document}

\begin{abstract}
    
\end{abstract}

\maketitle

% INTRO FOR NOW
This paper presents an Alloy model~\cite{alloy} for an AI agent
program where the agent can read and write to files, and create pull
requests (PRs). The goal of this paper is to show the benefits
of formal methods (FM) applied to agentic programs. This paper
is divided the following way:

% PROBLEM
A problem one can highlight is the separation between the purpose
of the AI Agent and the tools it has available. Tools are concrete
abstractions of what the agent is allowed to do, while a purpose
is a natural-language description of the goal. Here lies the problem:
Natural-language is ambiguous. We redefine purpose to the less ambiguous
noun operation, more specifically an agent has a set of operations. 
We define an operation as a specific task that lies within the allowable
actions a tool has. We can summarize in a more general way:
Agentic programs rely on natural-language leading to ambiguous requirements.

There are news, articles, blogs, etc. about the problems created by
under specified purposes.

% POTENTIAL SOLUTION
FM can be a potential solution by defining a more strict language
that removes the ambiguiuty of natural-languages.
\bibliography{refs}
\bibliographystyle{plain}

\end{document}